\begin{abstract}
Surgical training needs objective assessment of suturing skill, and rating
tools such as EASE assess how a surgeon holds the needle against defined
criteria. Rating every stitch by hand is costly, and automatic scoring so far
predicts the score end to end, as a black box. We propose to
measure rather than predict: EASE defines its criteria geometrically, and we
measure what they define. Detected events divide a monocular suturing video
into stitches and into the window that each criterion refers to. In that
window, keypoints and needle masks from a general-purpose vision--language
model, given a few examples, yield the grasp position, the hold angle, and
the hold depth in 3D, which the EASE criteria turn into scores. Every score
is thus traceable to the frames and the value behind it. On JIGSAWS, running
fully automatically from the raw video to the scores, the method reproduces a
surgeon's EASE scores with AUC 0.89--0.97, at least as well as direct grading
by a state-of-the-art VLM.
\end{abstract}
%
\begin{keywords}
Surgical skill assessment, surgical video, event detection, geometric
measurement
\end{keywords}
